% =====================================================================
% Study-domain overview figure.
%
% Source: projects/tessera_downscaling/scripts/maps/plot_region_overview.py
% Artefacts: scripts/maps/outputs/overview/
%   region_overview.pdf          the figure (both variables, one map)
%   region_station_counts.csv    per-region train / held-out / total
%
% Placement: works either as the first figure of Section 2.1 (it makes
% the "five climatically diverse regions ... chosen to span a wide range
% of terrain, land cover, and station density" sentence concrete before
% Table 1 quantifies it), or in Appendix A next to Table 3, whose bounds
% it draws. If it goes in the appendix, add a \Cref{} pointer from the
% "We evaluate across five climatically diverse regions" paragraph.
% =====================================================================

\begin{figure}[t]
  \centering
  \includegraphics[width=\textwidth]{figures/region_overview.pdf}
  \caption{%
    The five evaluation regions and the station network inside them.
    Black rectangles are the ERA5 crop boxes of \Cref{tab:region-boxes},
    which also serve as the region assignment for stations: the boxes do
    not overlap, and stations outside all of them are discarded, so every
    station belongs to exactly one region. Dots are the GHCNh stations
    that pass the selection of \Cref{sec:station-set}, for
    2\,m~temperature, 10\,m~wind~speed, or both, coloured by the fixed
    spatial split --- blue stations enter training, orange stations are
    withheld from training entirely and carry every score reported in
    this paper. The held-out 15\,\% is drawn at random, so it is
    interleaved with the training network rather than occupying a
    contiguous sub-domain. Region labels give the total station count $n$
    as \emph{2\,m~temperature\,/\,10\,m~wind~speed}, matching
    \Cref{tab:main}; the two networks overlap heavily but are not
    identical, since a station is counted per variable only where it
    reports that variable. Note the more than order-of-magnitude spread
    in network density across regions, from Europe to Australia.%
  }
  \label{fig:region-overview}
\end{figure}
